\section{One physical interval and four passage regions}

The preceding constructions use different phase coordinates. We now show
that their domains cover one connected interval of actual through
trajectories, with open overlaps. Their derivative estimates can then
be combined for the same affine displacement.

\begin{lemma}[A common entrance and gate choice]
\label{p2:lem-joint-entrance}
One auxiliary choice for the four passage regions, followed by
independent small coefficient bounds, has the following properties.
The incoming raw height $h$ is $x_{\rm raw}$ in
\eqref{deg:fixed-exit-natural-input} and $x_{\rm in}$ in
Lemma~\ref{lem:actual-dulac-matching}; its stopped value is $x_b$.
Thus the canonical weak input is exactly
\begin{equation}\label{p2:joint-raw-input}
 P=q(h)=\frac{h-x_b}{s_{\rm in}},\qquad
 s_{\rm in}=\Omega a^{2-\lambda_\gamma}U_{\rm in}>0,
 \qquad \Omega=r_s^{2-\lambda_\gamma},\quad U_{\rm in}^{\pm1}\le C.
\end{equation}
This is the affine coordinate on the original raw section, distinct
from the displacement coordinate $x=h-a_\theta$; no new input
observer is inserted.

Let $\widehat p_0=p_0$ be a fixed buffered raw cutoff in
\eqref{deg:fixed-exit-middle-invariant}, and choose
$0<\widehat p_*<\widehat p_0$ so that
$\mathfrak j'=(\Phi_\theta/\Phi_\theta')'>3/4$ on
$[0,\widehat p_*]$. Set $q_{\rm ref}=\widehat p_*/4$ and use
$P_\dagger=q_{\rm ref}$ in
\eqref{deg:source-growing-positive-log-phase}. The final negative
thresholds are denoted by
$\widehat r_{\rm ge}^{\rm out}=r_{\rm ge}^{\rm out}$ and
$\widehat r_{\rm ge}=r_{\rm ge}$, with a fixed positive gap. Every
actual first-lip input with $r_\ell\ge r_{\rm ge}$ satisfies
\begin{equation}\label{p2:joint-entrance-buffer}
 0<P<q_{\rm ref}/2<q_{\rm ref}<\widehat p_*<\widehat p_0,
\end{equation}
including first hits outside the growing product range.

The source constant $\delta_\sharp(L)$ is selected once for the full
bounded-source domain and relative negative-curvature estimate. Its gates
are those below in Lemma~\ref{p2:lem-bounded-source-budget}.
The final $r_{\rm ge}$ precedes $R_{\rm BD}$ and its annular
overlap; all fixed gates and phase buffers precede the final coefficient
bounds. No comparison of the independent scales is imposed.
\end{lemma}
\begin{proof}
Fix the original raw windows and a finite $V_1>1$, say $V_1=2$.
Choose the local cuts and ordinary $C^3$ section maps for this value in
Lemma~\ref{lem:actual-dulac-matching}, retaining its enlarged
interval $0\le V\le V_1+1$. Keep these choices throughout.
Use the strict source block and the two comparison orbits of
Proposition~\ref{bd:prop-bounded-action}: fix its derivative constants,
then $B_-$, then $\delta_\sharp(L)$ for both the domain and relative
estimate, and finally $\delta_1,\delta_2$ and the gates below.
Fix $r_c>2$ and choose $\ell_0$ below the ratio of a fixed real
comparison orbit at $-r_c$, as in
\eqref{deg:degree-ratio-monotonicity}.

Fix a source envelope $\bar\sigma_{\rm src}=1/3$ for the separate
centre-map estimates on $r_c\le r\le\bar\sigma_{\rm src}/s_0$ in
Lemma~\ref{deg:lem-growing-matching-direct}; its derivative constants
are fixed before the positive cutoff and negative thresholds below.
The growing and output intervals will be chosen inside this envelope.
The independent stopped outer chord and local invariant give the
positive input map on $[0,p_0]$ in
\eqref{deg:fixed-exit-natural-input}--\eqref{deg:fixed-exit-exact-weak-coordinate}.
Choose the smaller fixed $\widehat p_*$ above so that the positive
action at $P_0=q_{\rm ref}/2$ also exceeds the fixed middle
variational exponent by two. This is possible since
$\Phi_\theta(P)=PU_\theta(P)$ with uniform positive units and
$\tau/\chi$ a positive unit. These choices precede the final negative
threshold.

The full first-hit cutoff following \eqref{deg:degree-phase-order}
provides a fixed $R(P_0)$. Choose $R=r_{\rm ge}^{\rm out}$ larger
than this and every earlier growing threshold, and so that
$c(R-r_c)>C_*(P_0)+2d_*$ for a fixed $d_*>0$ in
\eqref{deg:source-growing-positive-endpoint-action-budget}.
Choose $r_{\rm ge}>R$ with a fixed positive buffer.
The separate centre maps give
\eqref{deg:source-growing-positive-image-inclusion} with endpoint
$P_0$ instead of $p_*$ after independent coefficient shrink.
Their monotonicity and scalar uniqueness therefore construct the
actual product reference $r_\ell=r_{\rm ge}$ with raw input
$0<P_{\rm ge}<P_0$, in the original buffered chord. The local
invariant supplies its positive-time continuation to the raw section,
as in \eqref{deg:source-growing-positive-raw-input-scale}.

On the full first-hit domain of the original through interval,
\eqref{deg:degree-phase-order} gives $r_{\ell,h}<0$.
Thus $r_\ell(h)\ge r_{\rm ge}$ implies
$h\le h_{\rm ge}=x_b+s_{\rm in}P_{\rm ge}$, and the same positive
affine input gives $0<P(h)\le P_{\rm ge}<P_0$.
This proves \eqref{p2:joint-entrance-buffer} throughout the first-hit
domain, including inputs beyond the product range. The extension follows
from order on the full domain; the product derivative estimates remain
restricted to their original range.

Only now choose $R_{\rm BD}$ so the compact strong projection of
$r_\ell=r_{\rm ge}$ lies strictly below $2R_{\rm BD}$.
The transverse inverse and first-hit order give a fixed
$r_{\rm ge}^{\sharp}>r_{\rm ge}$ in the enlarged annular overlap.
The annular proof now fixes its own outer cuts \(L_{\rm a}\), after
this output bound and its limiting-action envelope. The earlier source
cut \(L\), \(h_L\), \(\delta_\sharp(L)\), and literal gates stay fixed;
both passages continue to the same raw ports. The original weak cuts
satisfy \(c_+=p_{\rm a}/8<p_-=p_{\rm a}/2\) and
\(c_-=p_{\rm a}/8<r_{\rm a}/2\), so the annular root-gap construction
uses this same auxiliary tuple. Only final independent coefficient bounds
are reduced to admit the new positive-height tail cuts.
Choose $0<\sigma_1<\sigma_2<\sigma_0<1/4$, the common
negative section $u_c$, its larger
$\sigma_0<\bar\sigma\le\bar\sigma_{\rm src}$ buffer, with
$\bar\sigma<|u_c|/2$, and strict
action gap $\Delta$ as in the construction after
\eqref{deg:degree-phase-order}. After these fixed choices, shrink
the independent coefficients so the fixed annular overlap lies in
the counted product range and $s_0r_{\rm ge}^{\sharp}<\sigma_1$.
All source, raw and shooting buffers retain strict margins.
All these restrictions are independent upper bounds on the coefficients.
After they have been imposed, the finite multiplier will be chosen.
\end{proof}

\begin{proposition}[One actual carrier for the four passage regions]
\label{p2:prop-unified-carrier-cover}
Use the common choice of Lemma~\ref{p2:lem-joint-entrance}, the original field,
through word, sections, raw windows, $L,h_L$ and closing maps.
The source constant $\delta_\sharp(L)$ is the once-selected constant
for the full original bounded-source domain and its relative sign;
$\delta_1,\delta_2$ keep their literal values from that tuple.
Let $I_\theta$ be the connected interval of actual through inputs in
this word, within the original buffered windows and positive graph
denominator domain, with
\[
             0<H_\theta(h):=z(\gamma_s;h)<z_{\rm cap}.
\]
Here $z_{\rm cap}$ is the original common fixed raw-height cap,
chosen within the standing high-height bound. The stopped trace $x_b$
is excluded. Write $h^0$ for the actual trace of $h_L=0$ and
$I_\theta^- = I_\theta\cap\{h\le h^0\}$.

\begin{enumerate}
\item The original through interval $(x_b,h^0)$ has the strictly increasing
shooting coordinate $W_\theta(h)=w(u_c;h)$ with image
$(0,W_\theta(h^0))$, where $W_\theta(h^0)\to u_c^2/2$ uniformly.
Thus every point of the full original enlarged chord
$\overline{\mathcal J_c^{\sharp\sharp}}$ corresponds to an actual through orbit,
with a strict height buffer. Its backward shot and both raw continuations
are that same orbit.
\item Let $h_S$ be the trace of $h_L=-\delta_2$ and let $h_A,h_B$
be the gate traces of $r_b=r_a/2,2R_{\rm BD}$, respectively. Put
\[
 B_\theta=I_\theta^-\cap([h_S,h^0]\cup[h_B,h_A]),\qquad
 h_T=\min\{h_S,h_B\}.
 \tag{PCov1}
\]
Let $I_{\ell,\theta}$ be the original first-ratio-hit domain. On it put
$\sigma(h)=s_0r_\ell(h)$. Wherever $r_\ell\ge r_{\rm ge}$ use
the original raw middle input
\[
 V(h)=\frac{w(u_+;h)}{\ell_0h_{\rm ax}(u_+)\rho(u_+)}
     =\frac{y(u_+;h)}{\varepsilon_+}>0,
 \quad h_{\rm ax}=-\delta p_d,\quad
 \varepsilon_+=\frac{\ell_0h_{\rm ax}(u_-)}{a^2J_-}.
 \tag{PCov2}
\]
All factors in this definition are actual and positive for a physical
coefficient fibre. Define the following subsets of that same carrier:
\[
\begin{aligned}
 G_\theta&=I_\theta^-\cap I_{\ell,\theta}
          \cap\{r_\ell\ge r_{\rm ge},\ \sigma\le\sigma_0\},\\
 D_\theta&=I_\theta^-\cap I_{\ell,\theta}
          \cap\{\sigma\ge\sigma_1,\ 0<V\le V_1\},\\
 E_\theta&=I_\theta^-\cap I_{\ell,\theta}
          \cap\{\sigma\ge\sigma_1,\ V\ge1\}.
\end{aligned}                                                    \tag{PCov3}
\]
The $V$ coordinate in the last two lines is well defined because
$\sigma_1/s_0>r_{\rm ge}$. These are connected, relatively closed
intervals, possibly empty, and
\[
 I_\theta^-=B_\theta\cup G_\theta\cup D_\theta\cup E_\theta.
 \tag{PCov4}
\]
Each of $G_\theta,D_\theta,E_\theta$ has the exact raw input buffer
\[
 0<q(h):=\frac{h-x_b}{s_{\rm in}}<q_{\rm ref}
       <\widehat p_*<\widehat p_0,
 \qquad s_{\rm in}=\Omega a^{2-\lambda_\gamma}U_{\rm in}>0.
 \tag{PCov5}
\]
The compact annular/growing overlap and the original product and
$V$ overlaps are on these actual inputs, with their equality dividers
included. Intersecting with any original closed raw-window or height
truncation preserves the cover; no divider is required to occur.
\item $G_\theta$ lies in the growing product domain, and
$D_\theta,E_\theta$ lie in the enlarged domains of the finite-$V$ and
fixed-exit maps constructed in
Proposition~\ref{p2:prop-finite-v-fe-actual-budget}.
The high-height interval is $I_\theta^+=I_\theta\cap\{h\ge h^0\}$.
Together with (PCov4) it covers the whole $I_\theta$, including every
non-root through input. On all these intervals the physical equation is
\[
 F(h-a_\theta)=h-a_\theta-f_\theta(m\mathfrak T_\theta(h)),
 \qquad \mathfrak T_\theta=d_\theta\circ T_\theta^\infty.
 \tag{PCov6}
\]
\end{enumerate}
The statements about actual raw maps concern $j>0$; the reciprocal
face $j=0$ has only the uniform one-sided normalized interpretation.
\end{proposition}
\begin{proof}
The independently constructed stopped and $h_L=0$ graphs bound the
through interval. Every intermediate graph reaches $u_c\in(-1,u_-)$ with
positive denominator. Scalar first variation makes $W_\theta$ strictly
increasing, with stopped limit zero.
For the upper comparison graph set $r_0=s_0L$, $r=-u$ and $\mathscr W=w/r^2$.
At its initial cut
\[
 \mathscr W(r_0)=\tfrac12+O_L(\tau b_\tau),\qquad
 b_\tau=-\log(2\tau).
\]
On $r_0\le r\le |u_c|$, direct substitution of the original polynomial
in the graph equation gives, on $1/4\le\mathscr W\le3/4$,
\[
 r\mathscr W_r=1-2\mathscr W+\mathscr E,\qquad
 |\mathscr E|\le C\{\kappa\tau+\kappa\tau/r\}
                 \le C\{\kappa\tau+\tau/L\}.             \tag{PCov7}
\]
Indeed $|\delta p_d(-r)|\le C\kappa\tau(r+a)$ and
$a/r\le x_{\max}/L$, while $K_d=1+O(\kappa\tau)$ and
$-L_d(-r)/r=1+O(\kappa\tau)$; the $E_dw$ term is smaller on this
corridor. The regular reciprocal factorization gives the same bounds
uniformly as $j\downarrow0$. Both boundary vectors point inward.
The logarithmic-radius equation contracts, so variation of constants
gives $\mathscr W(|u_c|)\to1/2$. Hence
$W_\theta(h^0)\to u_c^2/2$. The uniform endpoint margins put
$\overline{\mathcal J_c^{\sharp\sharp}}$ strictly inside this shooting
image after shrink. Strict first variation gives a unique through input
for each $W$; backward uniqueness identifies its shot and the two
already constructed raw continuations. This proves (i) independently
of the fixed-exit phase.

The same incoming first variation preserves the ordering of the
actual $h_L$ labels. Thus the central source interval is $[h_S,h^0]$.
The exact gate inverse and centre matching give the annular interval
$[h_B,h_A]$. The fixed orbit $h_L=-\delta_1$ has trace $h_C$ strictly
inside both intervals, because $-\delta_2<-\delta_1<0$ and its actual
gate labels lie strictly inside the annular bands. Consequently
\[
 [h_S,h^0]\cup[h_B,h_A]=[h_T,h^0].                       \tag{PCov8}
\]
The first fixed comparison alone supplies this overlap. The compact overlap also supplies an actual reference with $r_\ell=r_{\rm ge}$ and $r_b<2R_{\rm BD}$;
its trace $h_{\rm ge}$ satisfies $h_B<h_{\rm ge}$.
Every input failing (PCov8) is below $h_T\le h_B$, so the supplied full
original first-hit order gives $r_\ell>r_{\rm ge}$. The unified entrance
then gives (PCov5) on all inputs with $r_\ell\ge r_{\rm ge}$, including
first hits beyond the product range. Their derivative estimates are applied only within the growing product
domain.

For this actual local input $q$, the original positive cut is reached
on the stopped outer chord's local invariant. Its height is positive,
finite and strictly increasing in $h$. Formula (PCov2) is exactly the
original normalization $w=\ell_0h_{\rm ax}\rho V$, so
$V_h>0$, while the full first-hit order gives $\sigma_h<0$.
All sets in (PCov3) are therefore intervals and include their finite
level boundaries. If an input outside $B_\theta$ has
$\sigma\le\sigma_0$, it is in $G_\theta$. Otherwise
$\sigma>\sigma_0>\sigma_1$ and, because $V_1>1$, its positive $V$
satisfies $V\le V_1$ or $V\ge1$. It is in $D_\theta$ or
$E_\theta$. This proves (PCov4) for every through input, whether or not $F$ vanishes.

The compact band $r_\ell\in[r_{\rm ge},r_{\rm ge}^{\sharp}]$
is strictly inside the annular interval and, after the prescribed
$s_0$ bound, the growing product domain. The inequalities
$\sigma_1<\sigma_2<\sigma_0\le\bar\sigma<|u_c|/2$ retain the actual
product overlap. At $\sigma=\sigma_0$ the other passage region has a strict
$\sigma_1$ buffer, and at $V=1,V_1$ the last two regions overlap
because $1<V_1$. If a level is absent, the corresponding truncated interval is empty. The bounded endpoints and annular/growing overlap are included in these
intervals.

On $G_\theta$, $r_\ell\ge r_{\rm ge}>r_{\rm ge}^{\rm out}$ and
$\sigma\le\sigma_0\le\bar\sigma$, so the independently constructed
negative lip and positive chord identify the actual growing graph.
On $D_\theta\cup E_\theta$ the same raw buffer (PCov5) supplies the
original positive local cuts. The construction in Proposition~\ref{p2:prop-finite-v-fe-actual-budget}, using (i) as its
independent shooting geometry, identifies the bounded-$V$ early-cut
passage and the fixed-exit shooting map with precisely these orbits. In particular its action bound places the whole $E_\theta$
inside the enlarged inverse image, with $\Theta=0$ buffered. This last step is the only use of that proposition; the shooting
geometry in part (i) was constructed independently.

Finally the $h^0$ comparison graph has original source height
$Y_0\asymp b_\tau\to\infty$ and physical height
$H_\theta(h^0)=s_s^2Y_0\to0$. The original source dictionary gives
$|\alpha_0|+\alpha_1^2=s_s^2(1+p_s^2)$, with $p_s^2=O(\tau)$.
It follows that
\[
 K(|\alpha_0|+\alpha_1^2)<H_\theta(h^0)<z_{\rm cap}.
\]
Scalar height order puts every retained input at or above $h^0$ in
that same original high-height interval. The real enlarged $h_L=0$
source chord supplies its open overlap with the low carrier. Thus
$I_\theta^+$ includes the equality $h^0$ and all its non-root inputs;
the upper cap is common and fixed, rather than selected at a zero.
Original-time uniqueness gives (PCov6), with no raw gain inverted in
the domain construction.
\end{proof}

\section{Bounded source}

\begin{lemma}[Physical derivative estimates on the bounded-source intervals]
\label{p2:lem-bounded-source-budget}
Keep the original field, $L,h_L$, raw ports and enlarged windows, regular
maps and ordinary $C^3$ divided closing. Keep the unified tuple's literal
\[
 \begin{gathered}
 \delta_1=\min\{\delta_\sharp(L)/8,x_{\min}^2/8192,1/8192\},\\
 \delta_2=2\delta_1,\quad p_{\rm cmp}=\sqrt{2\delta_1},\quad r_a=p_{\rm cmp}/8,\quad p_a=r_a/16,\\
 p_-=p_a/2,\quad c_\pm=p_a/8.
 \end{gathered}
\]
The constant $\delta_\sharp(L)$ was fixed once, sufficiently small for
the bounded-source domain and relative estimate, before these gate
constants were defined. On the actual source word put
\[
 \beta=h_L/\tau-\log(2\tau),\quad b_\tau=-\log(2\tau),\quad
 \beta_{\min}=-\delta_2/\tau-\log(2\tau).
\]
There are fixed thresholds $B_-,B_+$ and constants $C_B,C_H$, uniform in
the original compact and reciprocal coefficient families, such that:
\begin{enumerate}
\item On every actual input with $\beta_{\min}\le\beta\le-B_-$,
the complete divided transition $\Pi$ has $\Pi''<0$ in the original
affine raw input, hence $F''=-m\Pi''>0$ for $m>0$.
\item On $-B_-\le\beta\le B_+$, one has $0<\Pi'\le C_B$.
\item On $B_+\le\beta\le b_\tau$, including $h_L=0$, and on a real
enlarged source chord through that orbit, one has $0<\Pi'\le C_H$.
The same bound holds on its admitted height-dominant continuation.
\end{enumerate}
All assertions hold on enlarged overlap intervals and at every nonclosed
through input. One final choice $m\max\{C_B,C_H\}<1/2$ gives $F'>1/2$
on the last two pieces. The input is always
$x=h-a_\theta$, and
\[
 \Pi(x)=\mathcal C_{\theta,m}(T_\theta(a_\theta+x)),\qquad
 F(x)=x-m\Pi(x)=x-f_\theta(m d_\theta(T_\theta(a_\theta+x))).
 \tag{BS1}
\]
\end{lemma}

\begin{proof}
Apply Proposition~\ref{bd:prop-bounded-action} to the same sections
and divided closing. Its estimates
\eqref{bd:relative-negative-curvature}--\eqref{bd:relative-raw-curvature}
give the first assertion, and \eqref{bd:compact-raw-variation}
gives the second.
The independent frozen windows, $L_f=8$, $R_f=5$ and polynomial
margin $3/4$ are supplied by Lemma~\ref{bd:lem-source-ranges};
both exact raw gains and their common-power cancellation are retained
in (BS1). It remains to check the literal gates and the upper chord.

The final $\delta_\sharp(L)$ is chosen there for both the domain and
relative curvature estimate, after $B_-$ and before the displayed literal
$\delta_1,\delta_2$. Since $\delta_2\le\delta_\sharp/4$,
the central block has enlarged lower chords down to
$h_L=-3\delta_1>-\delta_\sharp$. Compact labels
$[-B_--1,B_++1]$ and positive labels beginning below $B_+$ give actual
overlaps by monotone incoming first variation. A final $\tau$ shrink
makes $\beta_{\min}<-B_-$ and $b_\tau>B_++2$.

For positive labels choose the signed high-height constants
$K,z_*,C_H$ first, then $B_+$ large. Equation~\eqref{bd:positive-height}
gives, through $\beta=b_\tau$ and the enlarged chord below,
\[
 Y_0>K(1+p_s^2),\qquad
 H=s_s^2Y_0\le C_L\tau^3\sigma^4j^2(1+b_\tau)<z_*.
 \tag{BS10}
\]
To realize the upper chord, freeze $B=b_\tau$ in the original
positive-energy construction. For $\beta/B\in[3/4,3/2]$, its incoming
energy $\beta/B+B^{-1}\log(2/B)$ stays strictly inside $[1/4,3]$.
The energy defect is $O_L(B^{-1})$ and its variation integral tends
uniformly to zero; the same first-contact proof and positive inverse
therefore construct this entire chord. It has
$Y_0=\beta-\log\beta+O_L(1)$ and
$\eta(\pm L)=L^2/2+O_L(\tau b_\tau)$ at the unchanged cuts.
The original quadratic corridor and raw-window margins continue both
sides. In particular $0\le h_L\le\tau b_\tau/2$ is a real enlargement
above the $h_L=0$ orbit.

By (BS10), the signed high-height bound used in
Proposition~\ref{bd:prop-bounded-action}, namely
Corollary~\ref{rew:cor-high-derivative-signed}, applies throughout
this chord and its continuation
$K(|\alpha_0|+\alpha_1^2)\le H<z_*$, including the equality face.
It gives the third assertion, with the height scale frozen at each base point.

The full transported output, including reference values, has a common
bound $M_{\rm out}$ on the original section windows. Choose $m$ last
with $m\le m_0$, $mM_{\rm out}<a_{\rm cl}/2$ and
$m\max(C_B,C_H)<1/2$. Then (BS1) gives $F'>1/2$ on the compact,
positive and high pieces, while the negative curvature gives $F''>0$
for $m>0$. The enlarged estimates include threshold equalities and
every nonclosed input. At $m=0$, $F(x)=x$; collapsed sections denote
only normalized one-sided limits.
\end{proof}

\section{Growing exit}

\paragraph{The growing displacement and its curvature.}
The action derivative of the matched map is known, but the finite
closing is evaluated at its full translated output. We first bound that
output and its first two action derivatives, then use the logarithmic
slope to prove convexity in the original affine input.
All phase derivatives fix every coefficient. Use the already constructed
joint auxiliary tuple, the actual first lip $r=r_\ell$, and its original
enlarged domain
\[
 \widehat r_{\rm ge}^{\rm out}\le r\le\bar\sigma/s_0,\qquad
 s_0=\kappa\sqrt B,\qquad 0<\sigma_0<\bar\sigma<1/2 .
 \tag{GC1}
\]
The counted interval has $r\ge\widehat r_{\rm ge}$ and $s_0r\le\sigma_0$.
Write the actual affine raw input as
\[
 h=x_b+s_{\rm in}P,\qquad
 s_{\rm in}=\Omega a^{2-\lambda_\gamma}U_{\rm in}>0,\qquad
 \Omega=r_s^{2-\lambda_\gamma},\qquad U_{\rm in}^{\pm1}\le C .
 \tag{GC2}
\]
The joint entrance construction supplies $0<P<\widehat p_*/2$.
Set
\[
 X=\Phi_\theta(P)/\Phi_\theta(q_{\rm ref}),\qquad
 \mathcal E=X\partial_X=\mathfrak j(P)\partial_P,\qquad
 \mathfrak j=\Phi_\theta/\Phi_\theta'>0,\qquad
 e_s=\tau/\chi ,
 \tag{GC3}
\]
with the denominator held fixed. Let
\[
 \mathcal P(r)=P_s(-r)/B_\tau(-r),\qquad
 \mathfrak d=\mathcal P/r,\qquad \mathcal D=\mathfrak d\partial_r .
\]
The polynomial estimates and the matched action derivative give
\[
 \begin{gathered}
 \mathcal P\asymp1+r,\quad \mathcal P'>0,\quad
 \mathfrak d^{\pm1}\le C,\quad |\mathcal D\log\mathfrak d|\le C,
 \quad e_s^{\pm1}\le C,\quad \mathfrak j'>3/4,\\
 -\mathcal Er=e_s\mathfrak d\,n,\qquad n=1+\zeta,\qquad
 n\ge1/2,\qquad |\zeta|+|\mathcal E\zeta|\le C\tau .
 \end{gathered}
 \tag{GC4}
\]
Here $\mathcal E$ acts on the actual graph, including the negative
variable in the complete quotient row.

Retain the full actual output
\[
 \begin{gathered}
 T(r)=\Xi(W_\theta(r))
      =t_\dagger-a_{\rm grow}\mathcal B(r),\qquad
 a_{\rm grow}=\Omega s_0^2U_{\rm out}>0,\qquad
 \mathcal B'=rG,\quad G>0,\\
 W_\theta'=-s_0^2r\Gamma_\theta,\qquad
 G=\Gamma_\theta\,\Xi'(W_\theta)/(\Omega U_{\rm out}) .
 \end{gathered}
 \tag{GC5}
\]
Thus $T$ includes its coefficient-dependent translation, original
moving endpoint, tails and regular output section map.
The actual output lies in a fixed compact physical section window $J_T$.
The first action derivative of the landing, composed with this same
$\Xi$, satisfies
\[
 \begin{split}
 |\mathcal D\log G|
 &\le C\left\{
 \frac{\tau}{\sqrt B\,r^2}
       \left(1+\left|\log\frac{\tau}{\sqrt B\,r}\right|\right)
       +s_0\kappa\tau+s_0^2\mathcal P(r)\right\}\\
 &\le C\{\tau(1+|\log\tau|)+s_0\}.
 \end{split}
 \tag{GC6}
\]
The ordinary finite closing has $f_\theta(0)=0$, $f_\theta'>c_f>0$,
and uniform ordinary $C^3$ bounds on a buffered interval containing
$mJ_T$ for $0\le m\le\bar m$. It is the original closing after the
output map $d_\theta$ already retained in $T$.
The physical affine displacement is exactly
\[
 F_\theta(h)=h-a_\theta-f_\theta(mT_\theta(h)),\qquad
 T_\theta(h)=T(r(P(h))).
 \tag{GC7}
\]

\begin{proposition}[Convexity and zero count in the growing range]
\label{p2:prop-growing-physical-count}
On the whole counted growing interval, there are common constants
$C_T,C_{\rm cl}$ and a multiplier bound $m_0>0$, selected after the
independent coefficient bounds, such that
\[
 |T|\le C_T,\qquad 0<\mathcal ET\le C_T,\qquad
 |\mathcal E^2T|\le C_T .
 \tag{GC8}
\]
For every actual through input in that interval, whether or not
$F_\theta=0$,
\[
 \begin{gathered}
 \left|\mathcal E^j\{f_\theta(mT)-f_\theta(0)\}\right|
       \le C_{\rm cl}m,\\
 \left|\mathcal E^j\{\log f_\theta'(mT)-\log f_\theta'(0)\}\right|
       \le C_{\rm cl}m,\qquad j=0,1,2,\quad 0\le m\le m_0 .
 \end{gathered}
 \tag{GC9}
\]
For $0<m\le m_0$, the complete physical displacement is strictly convex:
\[
 F_\theta''(h)>0\qquad\hbox{at every point of the growing interval}.
 \tag{GC10}
\]
Consequently it has at most one critical point, that point is simple
and upward, and it has at most two isolated zeros counted with
multiplicity. At $m=0$, $F_\theta'(h)=1$.
The constants are uniform on the full prescribed passive coefficient
family and both coefficient graphs, with the reciprocal estimates
uniform for physical $j>0$ approaching its one-sided boundary.
\end{proposition}

\begin{proof}
\emph{Bounds on the full output and its action derivatives.}
Put $L(r)=-T_r=a_{\rm grow}rG(r)>0$ and
\[
 d_r=\min\{1,(\widehat r_{\rm ge}
              -\widehat r_{\rm ge}^{\rm out})/2\}>0 .
\]
The original lower buffer places every $[r-d_r,r]$ in (GC1).
Since $\mathfrak d^{-1}$ is bounded, (GC6) gives
\[
 |\partial_r\log L|
 \le r^{-1}+\mathfrak d^{-1}|\mathcal D\log G|\le C_L .
\]
Let $D_T$ be the diameter of $J_T$. Using the actual output values,
\[
 D_T\ge T(r-d_r)-T(r)
      =\int_{r-d_r}^rL(v)\,dv
      \ge d_re^{-C_Ld_r}L(r).
 \tag{GC11}
\]
Thus $L\le C$ without dividing by $a_{\rm grow}$, $\Omega$ or a raw gain.
Equation (GC4) now gives
\[
 \mathcal ET=e_s\mathfrak d\,Ln>0,\qquad
 \mathcal E\log(\mathcal ET)
 =-e_sn\{\mathcal D\log\mathfrak d+\mathcal D\log L\}
       +\frac{\mathcal E\zeta}{n},\qquad
 \mathcal D\log L=\mathfrak d/r+\mathcal D\log G .
 \tag{GC12}
\]
Every term on the right is bounded. Hence $\mathcal ET\le C$ and
$|\mathcal E^2T|\le C$; $|T|$ is bounded by the fixed compact $J_T$.
The reference value of the output is included in this bound, so (GC8)
applies at the actual argument of the closing map.

\emph{Composition with the finite closing.}
The full translated $T$ and its first two action rows satisfy (GC8)
on the buffered actual image. The ordinary $C^3$ closing has positive
derivative units there. Lemma~\ref{phys:lem-c3-adapters}, specifically
\eqref{phys:eq-absolute-closing-two}, therefore gives (GC9) for all
through inputs, with constants fixed before $m$.

\emph{Strict convexity in the actual affine input.}
Before finite closing, the exact logarithmic slope is
\[
 \begin{split}
 H_0
 &:=\mathcal E\log
       \frac{\mathfrak j(P)}{e_sG(r)\mathcal P(r)n(P,r)}\\
 &=\mathfrak j'(P)+e_sn
       \left\{\frac{\mathcal P'(r)}r+\mathcal D\log G(r)\right\}
       -\frac{\mathcal E\zeta}{n}.
 \end{split}
 \tag{GC14}
\]
This identity follows from $\mathcal Er=-e_s\mathfrak d n$;
the derivative of the matched quotient is included.
By (GC4), (GC6) and $\mathcal P'>0$,
\[
 H_0\ge 3/4-C\{\tau(1+|\log\tau|)+s_0\}>1/2
 \tag{GC15}
\]
after common independent $\tau,\kappa$ bounds.
Here $s_0=\kappa\sqrt B$, with $B$ a positive unit.
The original passive coefficient set and joint auxiliary tuple are
unchanged.

Since $\mathcal ET=\mathfrak j T_P$ and (GC5) gives
$T_P=a_{\rm grow}e_sG\mathcal P n/\mathfrak j>0$, define
\[
 \mathscr R=\frac{s_{\rm in}}{mf_\theta'(mT)T_P}>0 .
\]
Every factor in its numerator, $a_{\rm grow}$ and $e_s$ is
coefficient-only. Its complete logarithmic row is therefore
\[
 H_m:=\mathcal E\log\mathscr R
    =H_0-m\frac{f_\theta''(mT)}{f_\theta'(mT)}\mathcal ET .
 \tag{GC16}
\]
The absolute output estimate (GC8) makes the closing term at most
$C_m m$, with a common $C_m$ fixed before the multiplier.
Choose last
$m_0\le\min\{\bar m,1,1/[4(1+C_m)]\}$. Then $H_m>1/4$.

Restore the physical input using (GC2). Direct differentiation of the
same displacement (GC7) gives at every through point
\[
 F_\theta'(h)=1-\frac1{\mathscr R},\qquad
 F_\theta''(h)=\frac{\mathscr R_h}{\mathscr R^2}
       =\frac{H_m}{s_{\rm in}\mathfrak j(P)\mathscr R}>0 .
 \tag{GC17}
\]
The two raw gains remain in this expression, so its sign holds even
when either gain is arbitrarily small.
Translation to $h-a_\theta$ leaves these derivatives unchanged.

Lemma~\ref{sharp:upper-critical-count} now gives the critical and
zero counts, with multiplicity; (GC17) excludes an identity interval.

The bound (GC6) uses $r\ge R_0$ and
$s_0^2\mathcal P(r)\le Cs_0$; its logarithm contains no $j$.
The bounds on $\kappa,\tau$ are independent, uniformly on the
prescribed compact $j$ ranges.
\end{proof}

\section{Finite height and fixed exit}
\subsection{The last two passage regions and their physical curvature}

Bounded amplified height is described by two small-power section maps.
For larger height we instead use the action coordinate on a fixed
negative section. We construct both maps on overlapping domains before
comparing their derivatives in the physical input.

Retain the single joint auxiliary choice, its original physical sections
and nested windows, the weak cuts, and the original coefficient family.
Let \(I_\ell\) be the positive-height first-lip domain in the original
through interval. Write
\[
 P={\frac{h-x_b}{s_{\rm in}}},\qquad
 s_{\rm in}=\Omega a^{2-\lambda_\gamma}U_{\rm in}>0,\qquad
 \Omega=r_s^{\,2-\lambda_\gamma},
 \qquad q_s={\frac{\beta}{\chi}}=\kappa Q_d ,
 \tag{FVFE1}
\]
where \(h\) is actual incoming height, \(Q_d,U_{\rm in}\) are positive
units, and no positive lower bound for either raw first gain is imposed.
The variable \(P\) is the canonical raw input of the joint entrance
construction; it is distinct from the saddle ratio \(q_s\).
Use the unchanged raw middle coordinate
\[
 V={\frac{y_+(h)}{\varepsilon_+}},\qquad
 \varepsilon_+={\frac{\ell_0h_-}{a^2J_m(u_-)}},\qquad
 \sigma=s_0r_\ell .
 \tag{FVFE2}
\]
Here \(y_+\) is the actual height \(w/a^2\) at \(u=u_+\).
Every phase derivative fixes all coefficients and all gains.

\begin{proposition}[Physical passage and convexity in the last two ranges]
\label{p2:prop-finite-v-fe-actual-budget}
For the joint auxiliary choice, followed by independent active
coefficient bounds, put
\[
 J_{34}=I_\ell\cap\{\sigma\ge\sigma_1\},\qquad
 D=J_{34}\cap\{0<V\le V_1\},\qquad
 E=J_{34}\cap\{V\ge1\},\qquad V_1>1 .
 \tag{FVFE3}
\]
These are intervals in the same actual incoming height, with every
displayed divider equality included. Each lies in the prescribed
enlarged domain of its corresponding actual transition, and
\(D\cup E=J_{34}\). Every input in these sets, whether
closed or nonclosed, has its original raw trace in
\(0<P<q_{\rm ref}<\widehat p_*<\widehat p_0\).
The finite-\(V\) map uses the original early negative cut on
the fixed enlarged interval \(0<V\le V_1+1\). For larger amplified
heights, the independently constructed fixed-exit map uses the buffered
shooting chord on \(u=u_c\). The two
constructions agree on \(1\le V\le V_1\) by original-time
uniqueness, so their union is the same physical transition on \(J_{34}\).

For the original return write
\[
 \mathcal F(h):=F(h-a_\theta)
       =h-a_\theta-f_\theta(mT(h)).
 \tag{FVFE4}
\]
The map \(T\) includes the actual negative output adapter. For every
\(0<m\le m_0\),
\[
 \mathcal F''(h)>0\qquad(h\in J_{34}).
 \tag{FVFE5}
\]
Thus the whole combined interval has at most one simple upward
critical point and at most two isolated zeros with multiplicity.
At $m=0$, $\mathcal F'=1$. The whole carrier is counted separately
in Proposition~\ref{p2:prop-physical-degree-count}.
\end{proposition}

\begin{proof}
\emph{The shooting chord and its phase coordinate.}
Part (i) of Proposition~\ref{p2:prop-unified-carrier-cover} realizes
$\overline{\mathcal J_c^{\sharp\sharp}}\Subset(0,u_c^2/2)$ as
unique through orbits, identified with the backward shots by uniqueness.
Proposition~\ref{fx:prop-inverse-phase} applies to this entire buffered
chord: its scalar first-exit construction is independent of the
positive raw-section dictionary. Its extracted phase
$\Psi_\theta(W)$ converges in $C^2$ to
$\Gamma_c\circ\mathcal W^{-1}$, whose derivative is
$1/H_{0,c}(r)>0$. Hence its increasing inverse
$W=\mathcal H_\theta(\Theta)$ has bounded first two derivatives
and a positive derivative unit on
\[
 \Theta_-\le\Theta\le\Gamma_- -\Delta,\qquad \Theta_-<0,
 \tag{FVFE8}
\]
strictly inside its inverse image. The negative buffer comes from the
original enlarged chord around $u_-$; the reciprocal functions give
the same uniform gaps.

\emph{Continuation of the shooting chord to the raw input.}
The exact positive invariant, including its local unit functions,
is
\[
 \Phi_\theta(P)=\{z_+y_+(P)\}^{\chi},\qquad
 \Phi_\theta(P)=P U_\theta(P),\qquad
 \Phi_\theta'>c,\qquad U_\theta^{\pm1}\le C.
 \tag{FVFE9}
\]
Use the reference $P_{\rm ref}=a_{\rm in}$ of
\eqref{deg:degree-fixed-exit-gap}:
\[
 \Phi_\theta(P_{\rm ref})=(z_+\beta\varepsilon_+)^\chi,\qquad
 q_s\log\Phi_\theta(P_{\rm ref})=-\Gamma_-+o(1).
 \tag{FVFE10}
\]
The exact identity \eqref{p2:reference-telescope} proves this limit
on both graphs: on the reciprocal graph use its prescribed $d,j$
in place of $\mathcal Rc,1/c$. All factors remain regular at $j=0$,
so its error bound is unchanged and no $\beta\log j$ occurs.
This reference corresponds to $V=\beta$. On the compact shooting chord the limiting
gap from \(\Gamma_-\) has a fixed positive minimum \(d_{\rm fx}\).
The scalar shooting estimate and (FVFE10) give, after shrink,
\[
 \beta\log V\le\Gamma_- -d_{\rm fx}/2,\qquad
 \beta\log y_+\le-d_{\rm fx}/3 .
 \tag{FVFE11}
\]
Since \(\beta\le C\chi\), this implies \(y_+<\chi^2\).
Thus (FVFE9) can be inverted on the original invariant chord, and
\[
 \Phi_\theta(P)\le C e^{-d_{\rm fx}/(3q_s)},\qquad
 0<P\le e^{-d_{\rm fx}/(4q_s)}<\widehat p_0
 \tag{FVFE12}
\]
after a further independent \(\kappa\) bound.
In particular the physical distance from the stopped trace satisfies
\(\sup(h-x_b)\le C e^{-d_{\rm fx}/(4q_s)}\to0\) on the
entire shooting chord, since \(s_{\rm in}\le C\).
The two local points have the same invariant in the same larger
coordinate box. Positivity and monotonicity along that invariant level
give the intervening orbit segment; the independently actual raw
continuation identifies its raw trace with the shot through orbit.
This places the full chord within the independently selected fixed cutoff
$\widehat p_0$, before any zero selection.

\emph{The bounded-\(V\) passage.}
Lemma~\ref{lem:actual-dulac-matching} independently continues the
entire larger interval $0\le V\le V_1+1$ to the early cut
$u_0=-1+\eta_0>u_-$. Equations~\eqref{eq:early-saddle-smallness}
and~\eqref{eq:early-middle-identity} give its positive invariant
at most $Ce^{-c_0/(2q_s)}$. The inversion used in (FVFE12) places the whole interval strictly
inside the raw cutoff.
The assigned set \(D\) is its restriction to \(0<V\le V_1\).
On the negative early cut, the lemma's exact invariant is
\[
 I_-(w)=\alpha_\theta R,\qquad
 R=\beta^{-\beta}C(v)^\beta,\qquad
 v=\widehat V=V+o_{C^3}(\chi).                 \tag{FVFE14}
\]
Its fixed cut choice puts the full $R\le1+o(1)$ interval strictly
inside the outer positive-height chord. At $R=0$ the output is the
stopped height $w_->0$. Thus both actual tails, ordinary $C^3$
observers and finite closing are on their original buffered domains,
with their exact values and positive first gains retained.

\emph{The domains of the two passage maps.}
The joint entrance supplies the actual product reference
\(r_\ell=\sigma_1/s_0\) after the active shrink.
Direct landing followed by the independent inverse shooting phase
gives
\[
 \Theta(h_{\sigma_1})\longrightarrow\Gamma_c(-\sigma_1),
 \qquad
 \Gamma_- -\Gamma_c(-\sigma)
      =\log{\frac{c+\sigma}{c(1-\sigma)}} .
 \tag{FVFE15}
\]
On the reciprocal chart the gap is
\(\log((1+j\sigma)/(1-\sigma))\), uniformly positive on the fixed
product buffer. The retained choice of \(\Delta\) and its strict
buffer imply
\[
 0<g_0\le\Theta(h_{\sigma_1})\le\Gamma_- -2\Delta .
 \tag{FVFE16}
\]
The actual orders \(r_{\ell,h}<0\) and \(V_h>0\) imply that every
member of \(E\) has \(0\le\Theta=\beta\log V\le\Theta(h_{\sigma_1})\).
Thus the entire fourth range lies in the domain (FVFE8) of the inverse
phase map, including inputs whose orbits are not closed. The equality \(V=1\) has the
negative-action buffer there. Since
\(\mathcal H_\theta([0,\Gamma_- -2\Delta])
\Subset\mathcal J_c^\sharp\Subset\mathcal J_c^{\sharp\sharp}\)
uniformly after shrink, its whole closed output-height subset has
a strict actual shooting buffer.
Decrease the independent coefficient bounds further so that
\(\beta\log(V_1+1)<g_0/2\). Then
\(V(h_{\sigma_1})=\exp(\Theta(h_{\sigma_1})/\beta)>V_1+1\),
so the third range retains the full buffered finite-\(V\) passage.
The two constructions identify the same raw inputs and outputs on \(1\le V\le V_1\) by original-time uniqueness.
No physical input is relabelled at \(V=1\).
Their actual \(\sigma\ge\sigma_1\) domain has
\(r_\ell\ge\sigma_1/s_0>r_{\rm ge}\), so the joint entrance's full
first-hit order supplies the sharper raw bound
\(P<q_{\rm ref}<\widehat p_*<\widehat p_0\).

\emph{Physical convexity for bounded \(V\).}
For \(m>0\), the actual-section construction gives exactly
\[
 \begin{gathered}
 h=x_b+P_*K(v),\qquad
 f_\theta(mT(h))=b_*+a_*H(v),\\
 K(v)=\Psi(v^\chi),\quad
 H(v)=D(\beta^{-\beta}C(v)^\beta),\qquad P_*,a_*>0 .
 \end{gathered}\tag{FVFE17}
\]
Here all regular tails and finite closing are already included.
Put \(e_C=vC'/C>0\). Lemma~\ref{lem:actual-dulac-matching}
supplies the stronger normalized estimates below, with \(\Sigma_+\)
the coefficient-only positive invariant scale in its construction:
\[
 \begin{gathered}
 e_C=1+O(\epsilon_\theta),\qquad
 \left|{\frac{ve_C'}{e_C}}\right|\le C\epsilon_\theta,\qquad
 |D''/D'|\le C,\\
 \sup_t|t\Psi''(t)/\Psi'(t)|\le C\Sigma_+,\qquad
 \epsilon_\theta\le C\beta^{-N}e^{-c_0/(3\beta)},\qquad
 \Sigma_+\le Ce^{-c_0/(2q_s)} .
 \end{gathered}
\]
The exact Wronskian row, with \(R=\beta^{-\beta}C(v)^\beta\), is
\[
 \begin{split}
 \mathcal Q(v)
 &=v\left({\frac{K''}{K'}}-{\frac{H''}{H'}}\right)\\
 &=\chi-\beta e_C-{\frac{ve_C'}{e_C}}
       +\chi v^\chi{\frac{\Psi''}{\Psi'}}
       -\beta e_C R{\frac{D''}{D'}}
 \ge\chi/2>0 .
 \end{split}
 \tag{FVFE18}
\]
On this interval \(R\le1+o(1)\). With
\(q_s=\beta/\chi=\kappa Q_d\longrightarrow0\),
the displayed estimates give the relative bound
\[
 \frac{\mathcal Q(v)}{\chi}
 \ge1-Cq_s-C\frac{\epsilon_\theta}{\chi}-C\Sigma_+
 \ge\frac12 .
\]
Indeed, \(\beta=q_s\chi\) implies
\[
 \frac{\epsilon_\theta}{\chi}
 \le Cq_s\beta^{-N-1}e^{-c_0/(3\beta)}\longrightarrow0.
\]
The input-section error has weight \(\chi\Sigma_+\), whereas
the bounded output and closing error has weight \(\beta\).
Thus independent bounds on \(\kappa,\tau\) give this sign for
every \(v>0\), without extending the small-power compositions
in ordinary \(C^3\) to the excluded stopped endpoint.
Since \(dh/dv=P_*K'>0\), differentiation of the exact physical
displacement, without multiplying it by a nonconstant factor, yields
\[
 \mathcal F''(h)=
 {\frac{a_*}{P_*^2}}{\frac{H'}{(K')^2}}{\frac{\mathcal Q(v)}{v}}>0 .
 \tag{FVFE19}
\]

\emph{Physical convexity on the fixed-exit interval.}
The assigned set \(E\) has
\(0\le\Theta\le\Theta(h_{\sigma_1})\), strictly inside the
larger buffered action interval (FVFE8).
Use the common actual map \eqref{deg:fixed-exit-actual-output},
$\Xi(W)=t_b+a_{\rm out}D^n(W)$, $a_{\rm out}=\Omega U_{\rm out}>0$,
on its whole buffered height chord. Its ordinary $C^3$ normalized
bounds and $0<\Omega\le C$ retain the full output in the original
raw and closing windows. For the fixed closing buffer $a_{\rm cl}>0$,
impose on the whole chord
\[
 m_0\left\{|t_b|+
 a_{\rm out}\sup_{\mathcal J_c^{\sharp\sharp}}|D^n|\right\}
 <a_{\rm cl}/2 .
 \tag{FVFE20a}
\]
The braces are uniformly bounded on the original output window,
so this selects $m$ independently of the first gain.

Put \(\widehat H=D^n\circ\mathcal H_\theta\). The exact raw formulas are
\[
 T(h)=t_b+a_{\rm out}\widehat H(\Theta(P)),\qquad
 \Theta(P)=\beta\log\beta+
       q_s\log{\frac{\Phi_\theta(P)}{\Phi_\theta(P_{\rm ref})}}.
 \tag{FVFE21}
\]
Writing \(p=h-x_b=s_{\rm in}P\), retain the raw gain in the two rows
\[
 0<pT_h=
 a_{\rm out}\widehat H' q_s{\frac{P\Phi_\theta'}{\Phi_\theta}}
       \le C a_{\rm out}q_s ,
 \tag{FVFE22}
\]
\[
 p{\frac{T_{hh}}{T_h}}
 =P{\frac{\Phi_\theta''}{\Phi_\theta'}}
       -P{\frac{\Phi_\theta'}{\Phi_\theta}}
       +q_sP{\frac{\Phi_\theta'}{\Phi_\theta}}
                    {\frac{\widehat H''}{\widehat H'}}
 =-1+O(P+q_s)\le-\tfrac12 .
 \tag{FVFE23}
\]
The ordinary $C^3$ closing has derivative units and bounded
pre-Schwarzian on the whole sampled interval. (FVFE22)--(FVFE23) give
\[
 p\left\{{\frac{T_{hh}}{T_h}}
       +m{\frac{f_\theta''(mT)}{f_\theta'(mT)}}T_h\right\}
 \le-\tfrac12+Cma_{\rm out}q_s<-\tfrac14 .
 \tag{FVFE24}
\]
Therefore
\[
 \mathcal F''(h)=
 -mf_\theta'(mT)T_h
 \left\{{\frac{T_{hh}}{T_h}}
       +m{\frac{f_\theta''(mT)}{f_\theta'(mT)}}T_h\right\}>0 .
 \tag{FVFE25}
\]
No lower bound for \(s_{\rm in}\), \(\Omega\), or \(T_h\) was used.

The positive input order makes \(J_{34}\) one interval.
The two convexity estimates concern the identical physical \(F\)
on an overlapping cover, so (FVFE5) holds on their union.
Lemma~\ref{sharp:upper-critical-count} gives the critical and zero
counts with multiplicity, also for every finite collection approaching
the excluded stopped boundary.
The endpoints at the phase boundaries lie inside the enlarged domains
of their respective passage maps.
At \(m=0\) the original completed displacement has derivative one.
The combined reference logarithms, normalized section maps and
first-cocycle powers above are regular at the reciprocal face; none introduces a
\(\beta\log j\) or \(\tau\log j\) restriction. Actual physical maps
are asserted for \(j>0\), with \(j=0\) a one-sided normalized limit.
\end{proof}

\section{Counting zeros on the complete degree-loss interval}

Each local calculation either gives strict convexity of the physical
displacement or excludes critical points by a positive derivative.
Because the local maps agree on overlapping intervals, the same
alternative holds at every point of the connected through interval.

\begin{proposition}[Physical counting on the degree-loss carrier]
\label{p2:prop-physical-degree-count}
Fix the compact positive-unit ranges of the degree-loss graph, or
the reciprocal graph with $0<j\le j_0$. Use the joint auxiliary
choice of the growing entrance and its actual gates. There are
independent active coefficient bounds, a common raw-height cap and
a last finite-closing multiplier bound with the following property.
On the resulting connected interval $I_\theta$ of actual through
inputs on the original section, the complete affine displacement
\[
 F_\theta(x)=x-f_\theta\bigl(m\mathfrak T_\theta(a_\theta+x)\bigr),
 \qquad x=h-a_\theta,
 \tag{PC1}
\]
has at most one critical point. Every critical point is a simple
upward crossing of $F_\theta'$. Consequently it has at most two
isolated zeros, counted with multiplicity, and no identity interval.
The interval $I_\theta$ consists of the same through word in the
original nested windows, with positive passage denominator and
$0<z(\gamma_s)<z_{\rm cap}$ for one fixed physical height cap.
The assertion holds at every actual input of $I_\theta$, without
assuming its orbit is closed. The original field, marked sections,
nested raw windows, $L,h_L$ and closing map define this same return.
Here $\mathfrak T_\theta=d_\theta\circ T_\theta^\infty$ is the
original transported output and $f_\theta=C_\theta-C_\theta(0)$.
Thus $f_\theta(0)=0$ and at $m=0$ the displacement is $F_\theta(x)=x$.
No physical passage at $j=0$ or at a stopped boundary is asserted.
\end{proposition}

\begin{proof}
Proposition~\ref{p2:prop-unified-carrier-cover} supplies the closed
low cover $B_\theta\cup G_\theta\cup D_\theta\cup E_\theta$ and
its overlapping high continuation on the connected actual interval
$I_\theta$. It includes every divider equality and every nonroot input;
only the stopped endpoint is excluded.
All restrictions have the same affine displacement (PC1).

On the annular part of $B_\theta$,
Lemma~\ref{bd:lem-annulus} gives $\Pi''<0$ after the ordinary
$C^3$ divided closing is included. Its principal margin is
$\lambda_1/(2x_{\max}R_{\rm BD})$; the independent active bounds
make the errors smaller than half this margin.
Lemma~\ref{p2:lem-bounded-source-budget} gives the same sign on
negative source inputs and $F'>1/2$ on the compact, positive and
high source pieces. Propositions~\ref{p2:prop-growing-physical-count}
and~\ref{p2:prop-finite-v-fe-actual-budget} give $F''>0$ on
$G_\theta$ and on the whole combined $D_\theta\cup E_\theta$.
These signs are in the original affine input and include both exact
raw gains and the translated output inside the closing map.

Take the finite intersection of the independent active bounds, then
one multiplier bound for all derivative estimates and the full absolute
output envelope in the original closing buffer. Every critical point
of (PC1) consequently has $F''>0$, including those on the dividers.
Lemma~\ref{sharp:upper-critical-count} gives at most one simple
upward critical point and the two-zero bound with multiplicity.
An identity interval would contain a critical point with $F''=0$.
Section~\ref{sec:common-cover-conclusion} assembles the coefficient
cover and Hausdorff capture.
\end{proof}
